% ============================================================
%  Chapter 3 — System Model and Problem Formulation
% ============================================================
\chapter{System Model and Problem Formulation}
\label{chap:sysmodel}

This chapter formally defines the system model that underlies the
DRL-LSTM-LBRO framework. It describes the three-tier network
architecture, the resource models for edge cloudlets and the cloud
fallback, the IoT task model, the M/M/$c$ queuing framework used to
compute latency, and the Markov Decision Process (MDP) formulation
of the scheduling problem. The chapter concludes by stating the
multi-objective optimisation problem that the framework is designed
to solve.

% ─────────────────────────────────────────────────────────
\section{Network Architecture}
\label{sec:network_arch}

\subsection{Three-Tier Topology}

The system comprises three tiers arranged in increasing distance
from IoT end devices:

\begin{enumerate}
  \item \textbf{IoT Device Tier.} A large population of
  heterogeneous IoT devices---sensors, actuators, cameras, and
  wearables---continuously generates computational tasks. Each
  device has negligible local processing capability and must
  offload every task to a higher tier.

  \item \textbf{Edge Cloudlet Tier.} A set of $N$ edge cloudlets
  are deployed at or near network access points. Each cloudlet $i
  \in \{1, \ldots, N\}$ is a multi-server node with finite CPU
  capacity, RAM, and a bounded task queue. Cloudlets are
  heterogeneous: they differ in processing speed, memory, number
  of servers, and queue depth.

  \item \textbf{Cloud Tier.} A single cloud data centre with
  effectively unlimited compute and memory serves as a fallback
  destination. It is reached through a WAN link that adds a fixed
  propagation delay to every offloaded task.
\end{enumerate}

\subsection{Communication Model}

IoT devices communicate with the nearest cloudlet over a Local
Area Network (LAN) link. The propagation delay from device to
cloudlet $i$ is denoted $d_i^{\text{prop}}$ (seconds), which
varies with physical distance and is treated as a known constant
for each cloudlet. The WAN propagation delay to the cloud is
$d_{\text{cloud}}^{\text{prop}}$, fixed at 50\,ms in this work.
Transmission delays are neglected because task payloads are small
relative to link bandwidth.

% ─────────────────────────────────────────────────────────
\section{Cloudlet and Cloud Resource Model}
\label{sec:resource_model}

\subsection{Cloudlet Parameters}

Each cloudlet $i$ is characterised by the following parameters,
summarised in Table~\ref{tab:cloudlet_params}:

\begin{itemize}
  \item $M_i$: processing speed in millions of instructions per
  second (MIPS), equal for all cloudlets at 8000\,MIPS in this
  work.
  \item $R_i$: RAM capacity in megabytes.
  \item $c_i$: number of parallel servers (CPU cores).
  \item $K_i$: maximum queue length (tasks). A task arriving when
  the queue is full is dropped immediately.
  \item $d_i^{\text{prop}}$: edge propagation delay in seconds.
\end{itemize}

\begin{table}[htbp]
\centering
\caption{Cloudlet Parameters for the 3-Cloudlet Topology}
\label{tab:cloudlet_params}
\renewcommand{\arraystretch}{1.15}
\begin{tabular}{lcccccc}
\toprule
\textbf{Cloudlet} & \textbf{MIPS} & \textbf{RAM (MB)} &
\textbf{Servers} & \textbf{Queue} &
\textbf{Prop.\ Delay (s)} \\
\midrule
Cloudlet 1 & 8000 & 16384 & 4 & 30 & 0.002 \\
Cloudlet 2 & 8000 & 8192  & 3 & 25 & 0.003 \\
Cloudlet 3 & 8000 & 4096  & 2 & 20 & 0.004 \\
\bottomrule
\end{tabular}
\end{table}

\subsection{Runtime Resource State}

At each time step $t$, the observable state of cloudlet $i$
includes:
\begin{itemize}
  \item $\text{cpu\_util}_i(t) \in [0,1]$: fraction of CPU
  capacity currently in use.
  \item $\text{ram\_util}_i(t) \in [0,1]$: fraction of RAM
  currently allocated.
  \item $q_i(t) \in [0, K_i]$: number of tasks currently in the
  queue.
\end{itemize}

These quantities are computed from the set of tasks currently
being served and waiting in the queue at each cloudlet.

\subsection{Cloud Resource Model}

The cloud is modelled as an M/M/$\infty$ server with infinite
capacity, meaning it never drops tasks. Service time at the cloud
is computed from the task's instruction count and a fixed cloud
processing speed of $M_{\text{cloud}} = 32{,}000$\,MIPS.
End-to-end latency for a cloud-assigned task is:
\begin{equation}
  l_{\text{cloud}} = d_{\text{cloud}}^{\text{prop}}
  + \frac{\text{MI}}{ M_{\text{cloud}} }
  \label{eq:cloud_latency}
\end{equation}
where MI is the task's instruction count in millions.

% ─────────────────────────────────────────────────────────
\section{IoT Task Model}
\label{sec:task_model}

\subsection{Task Characteristics}

Each IoT task $\tau$ arriving at time step $t$ is described by a
tuple $\langle \text{MI}, r_{\tau}, d_{\tau} \rangle$, where:
\begin{itemize}
  \item $\text{MI}$: computational demand in millions of
  instructions, drawn from a Poisson distribution with mean
  $\bar{\lambda}_{\text{MI}}$.
  \item $r_{\tau}$: RAM required for execution in MB, drawn from
  a uniform distribution over $[r_{\min}, r_{\max}]$.
  \item $d_{\tau}$: deadline in seconds by which the task must
  complete to avoid a deadline violation penalty.
\end{itemize}

One task arrives per time step. The task generator produces
workload statistics calibrated against the Google Borg cluster
trace~\cite{reiss2011google}, ensuring that the instruction-count
and resource-demand distributions reflect a realistic mix of
short and long jobs.

\subsection{Task Acceptance and Dropping}

A task $\tau$ assigned to cloudlet $i$ is accepted if and only if:
\begin{equation}
  q_i(t) < K_i \quad \text{and} \quad
  \text{ram\_util}_i(t) \cdot R_i + r_{\tau} \leq R_i
  \label{eq:accept_condition}
\end{equation}
If either condition fails, the task is dropped and a penalty of
$-P_{\text{drop}}$ is applied to the reward signal. The cloud
never drops tasks (its queue is unbounded and RAM is unlimited in
the model).

% ─────────────────────────────────────────────────────────
\section{Queuing Model}
\label{sec:queue_model}

\subsection{M/M/\textit{c} Framework}

Each edge cloudlet is modelled as an M/M/$c$ queue with a finite
waiting room of size $K_i$. Task arrivals follow a Poisson process
with rate $\lambda$ (tasks per step), and service times are
exponentially distributed with rate $\mu_i = M_i / \bar{\text{MI}}$
(tasks per step). The traffic intensity for cloudlet $i$ is:
\begin{equation}
  \rho_i = \frac{\lambda}{c_i \, \mu_i}
  \label{eq:traffic_intensity}
\end{equation}

\subsection{Waiting Time via Erlang-C}

When a task is accepted by cloudlet $i$ and all $c_i$ servers are
busy, it waits in the queue. The expected waiting time is
computed using the Erlang-C formula. Let
\begin{equation}
  C(c_i, \rho_i) =
  \frac{ \dfrac{(c_i \rho_i)^{c_i}}{c_i!\,(1-\rho_i)} }
       { \displaystyle\sum_{k=0}^{c_i-1}
         \frac{(c_i \rho_i)^k}{k!}
         + \frac{(c_i \rho_i)^{c_i}}{c_i!\,(1-\rho_i)} }
  \label{eq:erlang_c}
\end{equation}
be the probability that an arriving task must wait. The expected
queueing delay is:
\begin{equation}
  W_i = \frac{C(c_i, \rho_i)}{c_i \, \mu_i \,(1 - \rho_i)}
  \label{eq:queue_delay}
\end{equation}

\subsection{End-to-End Latency}

The total end-to-end latency for a task accepted by cloudlet $i$
is:
\begin{equation}
  l_i = d_i^{\text{prop}} + W_i + \frac{\text{MI}}{M_i}
  \label{eq:edge_latency}
\end{equation}
Only completed tasks contribute to the average latency metric;
dropped tasks are excluded.

% ─────────────────────────────────────────────────────────
\section{MDP Formulation}
\label{sec:mdp}

The task scheduling problem is formalised as a Markov Decision
Process (MDP) $\mathcal{M} = \langle \mathcal{S}, \mathcal{A},
\mathcal{P}, \mathcal{R}, \gamma \rangle$.

\subsection{State Space}

At each time step $t$, the agent observes a state vector
$s_t \in \mathcal{S} \subset \mathbb{R}^{6N+5}$. The state
is constructed from three groups of features:

\begin{enumerate}
  \item \textbf{Current cloudlet state} (3 values per cloudlet):
  CPU utilisation, RAM utilisation, and normalised queue occupancy
  for each of the $N$ cloudlets.

  \item \textbf{LSTM load predictions} (2 values per cloudlet):
  predicted next-step CPU and RAM utilisation from the per-cloudlet
  GRU-LSTM predictor.

  \item \textbf{Task features} (3 values): normalised MI demand,
  normalised RAM demand, and normalised deadline of the arriving
  task.

  \item \textbf{Global features} (2 values): a global load
  indicator (mean CPU utilisation across all cloudlets) and a
  Jain's Fairness Index (JFI) measuring current load balance.
\end{enumerate}

For the 3-cloudlet baseline, the state dimension is
$6 \times 3 + 5 = 23$. For the 6-cloudlet experiment it expands
to $6 \times 6 + 5 = 41$, with no other architectural change.

\subsection{Action Space}

The action $a_t \in \{0, 1, \ldots, N\}$ specifies the
destination for the arriving task: values $0$ through $N{-}1$
correspond to the $N$ edge cloudlets, and value $N$ corresponds
to the cloud fallback. For the 3-cloudlet topology,
$|\mathcal{A}| = 4$; for the 6-cloudlet topology,
$|\mathcal{A}| = 7$.

\subsection{Transition Dynamics}

The environment transitions deterministically given the chosen
action: the task is added to the selected destination's queue (if
accepted), resource utilisations are updated, and the next task
is generated. From the agent's perspective the dynamics are
partially stochastic because future task arrivals are unknown.

\subsection{Reward Function}

The scalar reward $r_t$ at step $t$ is a weighted combination
of six terms:

\begin{equation}
r_t = w_1 r_{\text{drop}} + w_2 r_{\text{latency}}
    + w_3 r_{\text{deadline}} + w_4 r_{\text{energy}}
    + w_5 r_{\text{balance}} + w_6 r_{\text{overload}}
\label{eq:reward}
\end{equation}

\begin{itemize}
  \item $r_{\text{drop}} = -1$ if the task is dropped, $0$
  otherwise.
  \item $r_{\text{latency}} = -l_i / l_{\max}$, a latency
  penalty normalised by the maximum acceptable latency.
  \item $r_{\text{deadline}} = -1$ if $l_i > d_\tau$, $0$
  otherwise.
  \item $r_{\text{energy}} = -e_i / e_{\max}$, an energy penalty
  proportional to estimated execution energy.
  \item $r_{\text{balance}} = \text{JFI}(t) - 1$, rewarding even
  distribution of load across cloudlets.
  \item $r_{\text{overload}} = -\mathbf{1}[\text{cpu\_util}_i > \theta]$,
  a penalty for routing to an already heavily loaded cloudlet,
  where $\theta = 0.62$ is the overload threshold.
\end{itemize}

The weights $\{w_k\}$ are tuned so that drop-rate reduction is
prioritised, with latency and energy as secondary objectives and
balance and overload prevention as regularisation terms.

\subsection{Discount Factor}

The discount factor $\gamma = 0.99$ encourages the agent to
consider long-horizon consequences of each placement decision,
particularly the downstream effect on queue build-up at the
selected cloudlet.

% ─────────────────────────────────────────────────────────
\section{Optimisation Objective}
\label{sec:objective}

Let $\pi : \mathcal{S} \rightarrow \mathcal{A}$ denote a
deterministic scheduling policy. The objective is to find a
policy $\pi^*$ that maximises the expected discounted cumulative
reward over an episode of $T$ steps:

\begin{equation}
\pi^* = \arg\max_{\pi} \;
\mathbb{E}_{\pi} \!\left[
  \sum_{t=0}^{T} \gamma^t \, r_t
\right]
\label{eq:objective}
\end{equation}

This objective jointly minimises the task drop rate, average
end-to-end latency, energy consumption, and load imbalance in a
principled manner through the reward decomposition in
Equation~(\ref{eq:reward}). Because the reward is a scalar
combination of all criteria, the agent learns to navigate their
trade-offs rather than optimising any single one in isolation.
